\documentclass[10pt]{article}
\usepackage[a4paper,margin=0.55in]{geometry}
\usepackage{amsmath,amssymb,mathtools}
\usepackage{xcolor}
\usepackage{hyperref}
\usepackage{url}
\setlength{\jot}{5pt}
\numberwithin{equation}{section}
\allowdisplaybreaks

\title{$N_{2}=N_{2a}+N_{2b}$ (Group I, Row III): UV logs, plus prescription, and the combined simplified result}
\date{}

\begin{document}
\maketitle

\section{Check of the split $N_{2}=N_{2a}+N_{2b}$}
The colour factor of $N_{2}$ splits as
\begin{align}
&-f^{adc}f^{fbe}U^{db}(x)U^{ce}(z)+N_{c}U^{af}(y)+N_{c}U^{af}(x)-N_{c}U^{af}(x)\notag\\
&=\underbrace{\Big[-f^{adc}f^{fbe}U^{db}(x)U^{ce}(z)+N_{c}U^{af}(x)\Big]}_{N_{2a}}+\underbrace{\Big[N_{c}U^{af}(y)-N_{c}U^{af}(x)\Big]}_{N_{2b}},
\end{align}
and both parts keep the same kernel and prefactor, so the split is \textbf{correct}.

One typesetting slip: in your $N_{2a}$ the brace $\Bigg\{$ is opened once but closed twice, and the colour factor stands outside the $\int_{x,y,z,x',w'}$ although it depends on $x,y,z$. The consistent form is
\begin{align}
    \frac{d^{3}N_{2a\,\rm NLO}^{\rm}}{d^{3}k}  & =  \frac{1}{(2\pi)^{3}}\, \frac{g^{4}}{8\pi^{4}}\,\frac{1}{k^{+}}\,\int_{\Lambda}^{k^{+}-\Lambda}\, \frac{dp^{+}}{2\pi} \, \int_{x,y,z,x',w'}\,e^{-ik\cdot(w'-z)}\,e^{-ik\cdot (z-x)\,\frac{p^{+}}{k^{+}}}\,\frac{1}{\left[p^{+}\,(y-x)^{2}+(k^{+}-p^{+})\,(y-z)^{2}\right]} \notag \\ & \times\frac{(x'-w')^{i}\,(z-x)^{m}}{(x'-w')^{2}\,(z-x)^{2}}\Bigg[\frac{\delta^{im}\,d_{\perp}}{2\,(x-z)^{2}}\,\left[(y-x)^{2}-(y-z)^{2}\right]\,\frac{p^{+}\,(k^{+}-p^{+})}{k^{+\,2}}\notag\\ &-\frac{p^{+}}{(k^{+}-p^{+})}\,\left[\frac{(y-x)^{i}\,(x-z)^{m}}{(x-z)^{2}}-\frac{(y-x)^{i}\,(y-z)^{m}}{2\,(y-z)^{2}}\right]\notag \\ & +\frac{(k^{+}-p^{+})\,\delta_{im}}{k^{+}}\,\left[\frac{(y-z)\cdot (x-z)}{(x-z)^{2}}+\frac{(y-x)\cdot (y-z)}{2\,(y-x)^{2}}-\frac{\left[(y-x)^{2}-(y-z)^{2}\right]}{2\,(x-z)^{2}}\right]\notag\\ & -\frac{(k^{+}-p^{+})}{p^{+}}\,\left[\frac{(y-z)^{i}\,(x-z)^{m}}{(x-z)^{2}}+\frac{(y-x)^{m}\,(y-z)^{i}}{2\,(y-x)^{2}}\right]\notag \\ & +\frac{p^{+}\,\delta_{im}}{k^{+}}\,\left[\frac{(y-x)\cdot (x-z)}{(x-z)^{2}}-\frac{(y-x)\cdot (y-z)}{2\,(y-z)^{2}}-\frac{\left[(y-x)^{2}-(y-z)^{2}\right]}{2\,(x-z)^{2}}\right]\notag \\ &-\left[\frac{(y-x)^{m}\,(x-z)^{i}}{(x-z)^{2}}-\frac{(y-x)^{m}\,(y-z)^{i}}{2\,(y-z)^{2}}+\frac{(y-z)^{m}\,(x-z)^{i}}{(x-z)^{2}}+\frac{(y-x)^{i}\,(y-z)^{m}}{2\,(y-x)^{2}}\right]\Bigg]\notag \\ & \times\Big[U^{ab'}(x')-U^{ab'}(w')\Big]\,\Big[-f^{adc}\,f^{fbe}\,U^{db}(x)\,U^{ce}(z)+N_{c}\,U^{af}(x)\Big]\,\rho^{b'}(x')\,\rho^{f}(y).
\label{eq:N2a}
\end{align}
$N_{2b}$ is exactly as you wrote it: the same expression with colour $\big[N_{c}U^{af}(y)-N_{c}U^{af}(x)\big]$.

\section{Shorthand}
\begin{align}
&\mathcal{N}=\frac{1}{(2\pi)^{3}}\frac{g^{4}}{8\pi^{4}}\frac{1}{k^{+}},\qquad\frac{\mathcal{N}}{2\pi}=\frac{1}{(2\pi)^{3}}\frac{g^{4}}{16\pi^{5}}\frac{1}{k^{+}},\qquad \xi=\frac{p^{+}}{k^{+}},\qquad\lambda=\frac{\Lambda}{k^{+}},\\
&\Phi=e^{-ik\cdot(z-x)\frac{p^{+}}{k^{+}}},\qquad D=p^{+}(y-x)^{2}+(k^{+}-p^{+})(y-z)^{2},\qquad T^{im}=\frac{(x'-w')^{i}(z-x)^{m}}{(x'-w')^{2}(z-x)^{2}},\\
&B^{im}=\frac{(y-z)^{i}(x-z)^{m}}{(x-z)^{2}}+\frac{(y-x)^{m}(y-z)^{i}}{2(y-x)^{2}},\qquad A^{im}=\frac{(y-x)^{i}(x-z)^{m}}{(x-z)^{2}}-\frac{(y-x)^{i}(y-z)^{m}}{2(y-z)^{2}},\\
&\mathcal{C}^{(a)}=\Big[U^{ab'}(x')-U^{ab'}(w')\Big]\Big[-f^{adc}f^{fbe}U^{db}(x)U^{ce}(z)+N_{c}U^{af}(x)\Big]\rho^{b'}(x')\rho^{f}(y),\\
&\mathcal{C}^{(b)}=\Big[U^{ab'}(x')-U^{ab'}(w')\Big]\Big[N_{c}U^{af}(y)-N_{c}U^{af}(x)\Big]\rho^{b'}(x')\rho^{f}(y),\qquad
\mathcal{C}^{(b)}_{z}=\mathcal{C}^{(b)}\big|_{U(x)\to U(z)},\\
&K_{S}=\frac{(x'-w')\cdot(y-z)}{(x'-w')^{2}(y-z)^{2}}\,\frac{(y-x)\cdot(y-z)}{(y-x)^{2}(x-z)^{2}},\qquad P(\xi)=\xi(1-\xi)+\frac{\xi}{1-\xi}+\frac{1-\xi}{\xi}.
\end{align}

\part*{Part A: UV}

\section{Step 1: where the $x$ integral can diverge}
\begin{itemize}
\item \textbf{Large $x$} (for $0<p^{+}<k^{+}$): $T\sim1/|x|$, $D\sim p^{+}x^{2}$, and the brackets are $O(1)$, so the integrand falls like $|x|^{-3}$. Finite.
\item \textbf{$x\to y$}: the brackets contain at most $1/|y-x|$, which is integrable. At $p^{+}\to k^{+}$, $D\to k^{+}(y-x)^{2}$, but $\int_{0}^{k^{+}}dp^{+}/D=\log[(y-z)^{2}/(y-x)^{2}]/[(y-z)^{2}-(y-x)^{2}]$ is only a logarithm, which is integrable. Finite.
\item \textbf{$x\to z$}: terms of the form $1/(x-z)^{2}$, which diverge logarithmically in two dimensions. This is the \textbf{UV} region.
\end{itemize}

\section{Step 2: the kernel as $x\to z$ (common to $N_{2a}$ and $N_{2b}$)}
Put $r=x-z\to0$, $n=y-z$, $a=x'-w'$; then $y-x=n-r$ and $(y-x)^{2}-(y-z)^{2}=-2n\cdot r+r^{2}$. At $r=0$: $\Phi=1$ and $D=k^{+}n^{2}$. Keep the $1/r^{2}$ terms and average over directions with $\langle r^{i}r^{j}\rangle=\frac{r^{2}}{2}\delta^{ij}$ (in $d$ dimensions $\frac{r^{2}}{d}\delta^{ij}$ with $d_{\perp}=d$, which gives the same results):
\begin{align}
\text{$d_{\perp}$ term:}\quad&\frac{d_{\perp}}{2}\xi(1-\xi)\frac{(-a\cdot r)(-2n\cdot r)}{a^{2}r^{4}}\ \to\ \xi(1-\xi)\frac{a\cdot n}{a^{2}r^{2}},\\
\text{$(1-\xi)\delta$ term:}\quad&\text{bracket}=\frac{2n\cdot r}{r^{2}}+O(1)\ \Rightarrow\ -(1-\xi)\frac{a\cdot n}{a^{2}r^{2}},\\
\text{$-\tfrac{1-\xi}{\xi}B$ term:}\quad&T^{im}B^{im}=-\frac{a\cdot n}{a^{2}r^{2}}+O(1/r)\ \Rightarrow\ +\frac{1-\xi}{\xi}\frac{a\cdot n}{a^{2}r^{2}},\\
\text{$\xi\delta$ term:}\quad&\text{bracket}=\frac{2n\cdot r}{r^{2}}+O(1)\ \Rightarrow\ -\xi\frac{a\cdot n}{a^{2}r^{2}},\\
\text{last bracket:}\quad&\text{bracket}=-\frac{2n^{m}r^{i}}{r^{2}}+O(1)\ \Rightarrow\ +\frac{a\cdot n}{a^{2}r^{2}},\\
\text{$-\tfrac{\xi}{1-\xi}A$ term:}\quad&T^{im}A^{im}=-\frac{a\cdot n}{a^{2}r^{2}}+O(1/r)\ \Rightarrow\ +\frac{\xi}{1-\xi}\frac{a\cdot n}{a^{2}r^{2}}.
\end{align}
The constant terms cancel ($-(1-\xi)-\xi+1=0$). With $1/D\to1/(k^{+}n^{2})$:
\begin{equation}
\text{kernel}\ \xrightarrow{x\to z}\ \frac{1}{k^{+}}P(\xi)\,\frac{(x'-w')\cdot(y-z)}{(x'-w')^{2}(y-z)^{2}(x-z)^{2}}\ =\ \frac{1}{k^{+}}P(\xi)\,K_{S}\Big|_{x\to z}.
\label{eq:limit}
\end{equation}
This was checked numerically: the angle-averaged kernel divided by $K_{S}$ tends to $P(\xi)$ ($4.41$, $2.25$, $9.2011$ at $\xi=0.2,0.5,0.9$), and the remainder falls like $r^{2}$.

\section{Step 3: colour at $x\to z$}
Adjoint Wilson lines leave $f$ invariant, so $U^{db}(z)U^{ce}(z)f^{fbe}=f^{dcg}U^{gf}(z)$. Then
\begin{equation}
-f^{adc}f^{fbe}U^{db}(z)U^{ce}(z)+N_{c}U^{af}(z)=-f^{adc}f^{dcg}U^{gf}(z)+N_{c}U^{af}(z)=0,
\end{equation}
which was checked numerically for a random SU(3) matrix (residual $2\times10^{-15}$).
\begin{itemize}
\item \textbf{$N_{2a}$: UV finite.} $\mathcal{C}^{(a)}$ vanishes linearly as $x\to z$, so the integrand is at most $O(1/|x-z|)$, which is integrable. \emph{No UV log and no subtraction are needed for $N_{2a}$.}
\item \textbf{$N_{2b}$: UV divergent.} $\mathcal{C}^{(b)}\to\mathcal{C}^{(b)}_{z}\neq0$, so the coefficient of the divergence is \eqref{eq:limit}$\times\mathcal{C}^{(b)}_{z}$.
\end{itemize}

\section{Step 4: subtraction for $N_{2b}$}
Define $S=\frac{1}{k^{+}}P(\xi)\,K_{S}\,\mathcal{C}^{(b)}_{z}$. It has exactly the behaviour \eqref{eq:limit} at $x\to z$, and it is integrable at $x\to y$ and at large $x$ ($K_{S}\sim|x|^{-3}$). Write
\begin{equation}
N_{2b}=\underbrace{\mathcal{N}\!\int\!\frac{dp^{+}}{2\pi}\!\int e^{-ik\cdot(w'-z)}S}_{N_{2b\,\rm UV}}+\underbrace{\mathcal{N}\!\int\!\frac{dp^{+}}{2\pi}\!\int e^{-ik\cdot(w'-z)}\big[\text{integrand of }N_{2b}-S\big]}_{N_{2b\,\rm finite}}.
\end{equation}
The second term is UV finite, so $d=2$ ($d_{\perp}=2$) can be used in it.

\section{Step 5: the UV part $N_{2b\,\rm UV}$}
\textbf{5a. The $x$ integral in two dimensions (shows where it diverges).} Put $x=z+t$, $s=|t|$, $n=y-z$. Writing $n^{2}-ns\cos\varphi=\frac12[(n^{2}+s^{2}-2ns\cos\varphi)+(n^{2}-s^{2})]$ and using $\int_{0}^{2\pi}\frac{d\varphi}{A-B\cos\varphi}=\frac{2\pi}{\sqrt{A^{2}-B^{2}}}$:
\begin{equation}
\int_{0}^{2\pi}d\varphi\,\frac{n^{2}-ns\cos\varphi}{n^{2}+s^{2}-2ns\cos\varphi}=\pi+\pi\frac{n^{2}-s^{2}}{|n^{2}-s^{2}|}=\begin{cases}2\pi&s<|n|\\0&s>|n|\end{cases}
\ \Rightarrow\ \int_{x}\frac{(y-x)\cdot(y-z)}{(y-x)^{2}(x-z)^{2}}=2\pi\int_{0}^{|y-z|}\frac{ds}{s},
\end{equation}
which is divergent only at $s\to0$, i.e.\ $x\to z$ (UV).

\textbf{5b. In $d=2-\epsilon$.} Introduce a Feynman parameter and shift $t\to t+\alpha n$:
\begin{align}
\mu^{-\epsilon}\!\!\int d^{d}t\,\frac{(n-t)\cdot n}{(n-t)^{2}t^{2}}&=\mu^{-\epsilon}n^{2}\pi^{d/2}\Gamma\big(2-\tfrac d2\big)\int_{0}^{1}\!d\alpha\,(1-\alpha)\big[\alpha(1-\alpha)n^{2}\big]^{\frac d2-2}
=\pi^{d/2}\Gamma\big(2-\tfrac d2\big)\frac{\Gamma(\frac d2-1)\Gamma(\frac d2)}{\Gamma(d-1)}(\mu^{2}n^{2})^{-\frac\epsilon2}\notag\\
&=\pi\left[\frac{2}{\epsilon_{\rm UV}}+2\gamma_{E}+\log\frac{(y-z)^{2}\mu^{2}_{\overline{MS}}}{4}\right]+O(\epsilon),\qquad \mu^{2}_{\overline{MS}}=4\pi e^{-\gamma_{E}}\mu^{2},
\end{align}
which was checked with sympy.

\textbf{5c. The $p^{+}$ integral, done exactly} ($K_{S}\mathcal{C}^{(b)}_{z}$ does not depend on $p^{+}$):
\begin{align}
\int_{\lambda}^{1-\lambda}\!d\xi\Big[\xi(1-\xi)+\Big(\frac{1}{1-\xi}-1\Big)+\Big(\frac1\xi-1\Big)\Big]
&=\Big(\frac16-\lambda^{2}+\frac{2\lambda^{3}}{3}\Big)+2\log\frac{1-\lambda}{\lambda}-2(1-2\lambda)\notag\\
&=2\log\frac{k^{+}-\Lambda}{\Lambda}-\frac{11}{6}+\frac{4\Lambda}{k^{+}}-\frac{\Lambda^{2}}{k^{+2}}+\frac{2\Lambda^{3}}{3k^{+3}}.
\end{align}

\textbf{5d. Result.} The prefactor is $\frac{\mathcal{N}}{2\pi}\cdot\pi\cdot N_{c}$, so
\begin{align}
\frac{d^{3}N_{2b\,\rm UV\,NLO}}{d^{3}k}&=\frac{1}{(2\pi)^{3}}\frac{g^{4}N_{c}}{16\pi^{4}}\frac{1}{k^{+}}\left[2\log\Big(\frac{k^{+}-\Lambda}{\Lambda}\Big)-\frac{11}{6}+\frac{4\Lambda}{k^{+}}-\frac{\Lambda^{2}}{k^{+2}}+\frac{2\Lambda^{3}}{3k^{+3}}\right]\int_{y,z,x',w'}e^{-ik\cdot(w'-z)}\notag\\
&\times\left[\frac{2}{\epsilon}+2\gamma_{E}+\log\frac{(y-z)^{2}\mu^{2}_{\overline{MS}}}{4}\right]\frac{(x'-w')\cdot(y-z)}{(x'-w')^{2}(y-z)^{2}}\Big[U^{ab'}(x')-U^{ab'}(w')\Big]\Big[U^{af}(y)-U^{af}(z)\Big]\rho^{b'}(x')\rho^{f}(y)\notag\\
&\xrightarrow{\Lambda\to0}\ \frac{1}{(2\pi)^{3}}\frac{g^{4}N_{c}}{16\pi^{4}}\frac{1}{k^{+}}\left[{\color{blue}2\log\frac{\vee}{\Lambda}-2\log\frac{\vee}{k^{+}}}-\frac{11}{6}\right]\int_{y,z,x',w'}(\cdots),
\label{eq:UV}
\end{align}
using $2\log\frac{k^{+}}{\Lambda}=2\log\frac{\vee}{\Lambda}-2\log\frac{\vee}{k^{+}}$. This is the only place where $\log\mu$ appears.

\part*{Part B: plus prescription and the large longitudinal logs}

\section{Step 6: definitions}
\[
\Big[\frac{k^{+}}{p^{+}}\Big]_{+}f(p^{+})\equiv\frac{k^{+}}{p^{+}}\big[f(p^{+})-f(0)\big],\qquad
\Big[\frac{k^{+}}{q^{+}}\Big]_{+}h(q^{+})\equiv\frac{k^{+}}{q^{+}}\big[h(q^{+})-h(0)\big],\]\[
\Big[\frac{k^{+}}{k^{+}-p^{+}}\Big]_{+}f(p^{+})\equiv\frac{k^{+}}{k^{+}-p^{+}}\big[f(p^{+})-f(k^{+})\big],
\]
so that, exactly,
\begin{equation}
\int_{\Lambda}^{k^{+}-\Lambda}dp^{+}\frac{k^{+}}{p^{+}}f=\int_{0}^{k^{+}}dp^{+}\Big[\frac{k^{+}}{p^{+}}\Big]_{+}f+k^{+}f(0)\log\frac{k^{+}-\Lambda}{\Lambda}+O(\Lambda).
\end{equation}
The colour factors do not depend on $p^{+}$, so the same steps apply to $N_{2a}$ and $N_{2b}$.

\section{Step 7: the two singular terms of the kernel}
\textbf{7a. $p^{+}\to0$.} Write $-\frac{k^{+}-p^{+}}{p^{+}}B=-\frac{k^{+}}{p^{+}}B+B$ and $f=\Phi\,T^{im}B^{im}/D$. At $p^{+}=0$: $\Phi=1$ and $D=k^{+}(y-z)^{2}$, so
\begin{gather}
T^{im}B^{im}=\frac{(x'-w')\cdot(y-z)\,(x-z)\cdot(z-x)}{(x'-w')^{2}(z-x)^{2}(x-z)^{2}}+\frac{(x'-w')\cdot(y-z)\,(y-x)\cdot(z-x)}{2(x'-w')^{2}(z-x)^{2}(y-x)^{2}},\\
-k^{+}f(0)=-\frac{T^{im}B^{im}}{(y-z)^{2}}.
\end{gather}
\textbf{7b. $p^{+}\to k^{+}$.} Set $p^{+}=k^{+}-q^{+}$; then $\int_{\Lambda}^{k^{+}-\Lambda}dp^{+}=-\int_{k^{+}-\Lambda}^{\Lambda}dq^{+}$ and $-\frac{p^{+}}{k^{+}-p^{+}}A=-\big(\frac{k^{+}}{q^{+}}-1\big)A$, so
\begin{equation}
\int_{k^{+}-\Lambda}^{\Lambda}dq^{+}\Big(\frac{k^{+}}{q^{+}}-1\Big)h=\int_{k^{+}}^{0}dq^{+}\Big\{\Big[\frac{k^{+}}{q^{+}}\Big]_{+}-1\Big\}h-k^{+}h(0)\log\frac{k^{+}-\Lambda}{\Lambda},
\end{equation}
with $h=e^{-ik\cdot(z-x)\frac{k^{+}-q^{+}}{k^{+}}}T^{im}A^{im}/[(k^{+}-q^{+})(y-x)^{2}+q^{+}(y-z)^{2}]$. At $q^{+}=0$ the denominator is $k^{+}(y-x)^{2}$, and the phase combines with $e^{-ik\cdot(w'-z)}$ into $e^{-ik\cdot(w'-x)}$:
\begin{gather}
T^{im}A^{im}=\frac{(x'-w')\cdot(y-x)\,(x-z)\cdot(z-x)}{(x'-w')^{2}(z-x)^{2}(x-z)^{2}}-\frac{(x'-w')\cdot(y-x)\,(y-z)\cdot(z-x)}{2(x'-w')^{2}(z-x)^{2}(y-z)^{2}},\\ -k^{+}h(0)=-e^{-ik\cdot(z-x)}\frac{T^{im}A^{im}}{(y-x)^{2}}.
\end{gather}
Both endpoint coefficients were checked numerically: for random points and $\Lambda/k^{+}=10^{-7}$, the $\Lambda$-cut integral agrees with the plus form plus the logs to $10^{-8}$.

\textbf{7c. The subtraction of $N_{2b}$.} Write $P(\xi)=\xi(1-\xi)+\frac{k^{+}}{k^{+}-p^{+}}-1+\frac{k^{+}}{p^{+}}-1$. Each endpoint then gives $K_{S}\mathcal{C}^{(b)}_{z}\log\frac{k^{+}-\Lambda}{\Lambda}$:
\begin{equation}
-\frac{1}{k^{+}}\!\int_{\Lambda}^{k^{+}-\Lambda}\!\!dp^{+}P\,K_{S}\mathcal{C}^{(b)}_{z}=-\frac{1}{k^{+}}\!\int_{0}^{k^{+}}\!\!dp^{+}\Big[\xi(1-\xi)+\Big[\frac{k^{+}}{k^{+}-p^{+}}\Big]_{+}+\Big[\frac{k^{+}}{p^{+}}\Big]_{+}-2\Big]K_{S}\mathcal{C}^{(b)}_{z}-2\log\frac{k^{+}-\Lambda}{\Lambda}K_{S}\mathcal{C}^{(b)}_{z}.
\end{equation}
The first term is left unevaluated. Evaluated, it equals $+\frac{11}{6}K_{S}\mathcal{C}^{(b)}_{z}$.


\part*{Part C: adding $N_{2a}$ and $N_{2b}$ and simplifying}

\section{Step 8: adding the colour factors}
The kernels of $N_{2a}$ and $N_{2b}$ are identical, so their finite and log parts add simply by adding the colour factors:
\begin{equation}
\mathcal{C}^{(a)}+\mathcal{C}^{(b)}=\Big[U^{ab'}(x')-U^{ab'}(w')\Big]\Big[-f^{adc}f^{fbe}U^{db}(x)U^{ce}(z)+N_{c}U^{af}(y)\Big]\rho^{b'}(x')\rho^{f}(y)\equiv\mathcal{C}_{2}.
\end{equation}
The UV part and the subtraction belong to $N_{2b}$ only, and keep the colour $\mathcal{C}^{(b)}_{z}$.

\section{Step 9: simplifying the kernel}
Write $a=x'-w'$, $r=x-z$, $u=y-x$, $n=y-z$. Then
\begin{equation}
n=u+r,\qquad u^{2}-n^{2}=r^{2}-2\,n\cdot r,\qquad u\cdot r+n\cdot r=2\,n\cdot r-r^{2},\qquad n\cdot r-u\cdot r=r^{2}.
\end{equation}
Contract every bracket with $T^{im}=-a^{i}r^{m}/(a^{2}r^{2})$ and multiply by $a^{2}r^{2}$. Note $\delta^{im}T^{im}\,a^{2}r^{2}=-a\cdot r$.

\textbf{9a. $d_{\perp}$ term} ($d_{\perp}=2$ in the finite part):
\begin{equation}
\xi(1-\xi)\frac{u^{2}-n^{2}}{r^{2}}(-a\cdot r)=\xi(1-\xi)\,(a\cdot r)\,\frac{n^{2}-u^{2}}{r^{2}}.
\end{equation}

\textbf{9b. $(1-\xi)\delta$ term.} The bracket is
\begin{equation}
\frac{n\cdot r}{r^{2}}+\frac{u\cdot n}{2u^{2}}-\frac{u^{2}-n^{2}}{2r^{2}}=\frac{2\,n\cdot r}{r^{2}}-\frac12+\frac{u\cdot n}{2u^{2}},
\end{equation}
so this term gives $-(1-\xi)(a\cdot r)\big[\frac{2n\cdot r}{r^{2}}-\frac12+\frac{u\cdot n}{2u^{2}}\big]$.

\textbf{9c. $\xi\delta$ term.} The bracket is
\begin{equation}
\frac{u\cdot r}{r^{2}}-\frac{u\cdot n}{2n^{2}}-\frac{u^{2}-n^{2}}{2r^{2}}=\frac{2\,n\cdot r}{r^{2}}-\frac32-\frac{u\cdot n}{2n^{2}},
\end{equation}
so this term gives $-\xi(a\cdot r)\big[\frac{2n\cdot r}{r^{2}}-\frac32-\frac{u\cdot n}{2n^{2}}\big]$.

\textbf{9d. Last bracket:}
\begin{equation}
(a\cdot r)\frac{u\cdot r+n\cdot r}{r^{2}}-\frac{(a\cdot n)(u\cdot r)}{2n^{2}}+\frac{(a\cdot u)(n\cdot r)}{2u^{2}}
=(a\cdot r)\Big[\frac{2\,n\cdot r}{r^{2}}-1\Big]-\frac{(a\cdot n)(u\cdot r)}{2n^{2}}+\frac{(a\cdot u)(n\cdot r)}{2u^{2}}.
\end{equation}

\textbf{9e. Adding 9b, 9c and 9d.} The $\frac{2n\cdot r}{r^{2}}$ pieces cancel, because $-(1-\xi)-\xi+1=0$. The constants give $(a\cdot r)\big[\frac{1-\xi}{2}+\frac{3\xi}{2}-1\big]=(\xi-\frac12)(a\cdot r)$.

\textbf{9f. The $+B$ term} (left over from $-\frac{k^{+}-p^{+}}{p^{+}}B=-[\frac{k^{+}}{p^{+}}]_{+}B+B$), together with the two singular structures:
\begin{equation}
a^{2}r^{2}\,T^{im}B^{im}=-(a\cdot n)\Big[1+\frac{u\cdot r}{2u^{2}}\Big],\qquad
a^{2}r^{2}\,T^{im}A^{im}=-(a\cdot u)\Big[1-\frac{n\cdot r}{2n^{2}}\Big].
\end{equation}

\textbf{9g. Result.} The regular part of the kernel, including the $+B$ term, is
\begin{align}
\mathcal{R}(\xi)&=\xi(1-\xi)\,(a\cdot r)\,\frac{n^{2}-u^{2}}{r^{2}}+\Big(\xi-\frac12\Big)(a\cdot r)-a\cdot n
-(1-\xi)\frac{(a\cdot r)(u\cdot n)}{2u^{2}}+\xi\frac{(a\cdot r)(u\cdot n)}{2n^{2}}\notag\\
&\quad+\frac{(a\cdot u)(n\cdot r)-(a\cdot n)(u\cdot r)}{2u^{2}}-\frac{(a\cdot n)(u\cdot r)}{2n^{2}},
\end{align}
with the two endpoint structures
\begin{equation}
G_{p}=(a\cdot n)\Big[1+\frac{u\cdot r}{2u^{2}}\Big],\qquad G_{q}=(a\cdot u)\Big[1-\frac{n\cdot r}{2n^{2}}\Big],
\end{equation}
so that $a^{2}r^{2}\,T^{im}\big[R^{im}+B^{im}\big]=\mathcal{R}(\xi)$, $a^{2}r^{2}\,T^{im}B^{im}=-G_{p}$ and $a^{2}r^{2}\,T^{im}A^{im}=-G_{q}$. These three identities were checked numerically at random points (agreement to $10^{-15}$). In the original variables:
\begin{align}
\mathcal{R}(\xi)&=\xi(1-\xi)\,\frac{(x'-w')\cdot(x-z)\,\big[(y-z)^{2}-(y-x)^{2}\big]}{(x-z)^{2}}+\Big(\xi-\frac12\Big)(x'-w')\cdot(x-z)-(x'-w')\cdot(y-z)\notag\\
&-(1-\xi)\frac{(x'-w')\cdot(x-z)\,(y-x)\cdot(y-z)}{2(y-x)^{2}}+\xi\,\frac{(x'-w')\cdot(x-z)\,(y-x)\cdot(y-z)}{2(y-z)^{2}}\notag\\
&+\frac{(x'-w')\cdot(y-x)\,(y-z)\cdot(x-z)-(x'-w')\cdot(y-z)\,(y-x)\cdot(x-z)}{2(y-x)^{2}}-\frac{(x'-w')\cdot(y-z)\,(y-x)\cdot(x-z)}{2(y-z)^{2}},\\
G_{p}&=(x'-w')\cdot(y-z)\Big[1+\frac{(y-x)\cdot(x-z)}{2(y-x)^{2}}\Big],\qquad
G_{q}=(x'-w')\cdot(y-x)\Big[1-\frac{(y-z)\cdot(x-z)}{2(y-z)^{2}}\Big].
\end{align}

\section{Final result for $N_{2}=N_{2a}+N_{2b}$}
\begin{equation}
\frac{d^{3}N_{2\,\rm NLO}}{d^{3}k}=\frac{d^{3}N_{2\,\rm UV}}{d^{3}k}+\frac{d^{3}N_{2\,\rm finite}}{d^{3}k}+\frac{d^{3}N_{2\,\log}}{d^{3}k}+O\Big(\frac{\Lambda}{k^{+}}\Big),
\end{equation}
with the following three pieces.

\textbf{UV part} (from $N_{2b}$ only):
\begin{align}
\frac{d^{3}N_{2\,\rm UV}}{d^{3}k}&=\frac{1}{(2\pi)^{3}}\frac{g^{4}N_{c}}{16\pi^{4}}\frac{1}{k^{+}}\left[{\color{blue}2\log\frac{\vee}{\Lambda}-2\log\frac{\vee}{k^{+}}}-\frac{11}{6}\right]\int_{y,z,x',w'}e^{-ik\cdot(w'-z)}\notag\\&\times\left[\frac{2}{\epsilon}+2\gamma_{E}+\log\frac{(y-z)^{2}\mu^{2}_{\overline{MS}}}{4}\right]\frac{(x'-w')\cdot(y-z)}{(x'-w')^{2}(y-z)^{2}}\notag\\
&\times\Big[U^{ab'}(x')-U^{ab'}(w')\Big]\Big[U^{af}(y)-U^{af}(z)\Big]\rho^{b'}(x')\rho^{f}(y).
\label{eq:finUV}
\end{align}

\textbf{Finite part} ($d_{\perp}=2$; no $\Lambda$, no $\log\mu$):
\begin{align}
\frac{d^{3}N_{2\,\rm finite}}{d^{3}k}&=\frac{1}{(2\pi)^{3}}\frac{g^{4}}{8\pi^{4}}\frac{1}{k^{+}}\int_{0}^{k^{+}}\frac{dp^{+}}{2\pi}\int_{x,y,z,x',w'}e^{-ik\cdot(w'-z)}\Bigg\{\notag\\&\quad
e^{-ik\cdot(z-x)\frac{p^{+}}{k^{+}}}\,\frac{1}{\big[p^{+}(y-x)^{2}+(k^{+}-p^{+})(y-z)^{2}\big]}\,\frac{1}{(x'-w')^{2}(x-z)^{2}}\notag\\
&\qquad\times\Big[\mathcal{R}\Big(\frac{p^{+}}{k^{+}}\Big)+\Big[\frac{k^{+}}{p^{+}}\Big]_{+}G_{p}\Big]\Big[U^{ab'}(x')-U^{ab'}(w')\Big]\Big[-f^{adc}f^{fbe}U^{db}(x)U^{ce}(z)+N_{c}U^{af}(y)\Big]\rho^{b'}(x')\rho^{f}(y)\notag\\
&\quad-\frac{1}{k^{+}}\left[\frac{p^{+}(k^{+}-p^{+})}{k^{+2}}+\left[\frac{k^{+}}{k^{+}-p^{+}}\right]_{+}+\left[\frac{k^{+}}{p^{+}}\right]_{+}-2\right]\frac{(x'-w')\cdot(y-z)}{(x'-w')^{2}(y-z)^{2}}\,\frac{(y-x)\cdot(y-z)}{(y-x)^{2}(x-z)^{2}}\notag\\
&\qquad\times\Big[U^{ab'}(x')-U^{ab'}(w')\Big]\Big[N_{c}U^{af}(y)-N_{c}U^{af}(z)\Big]\rho^{b'}(x')\rho^{f}(y)\Bigg\}\notag\\
&+\frac{1}{(2\pi)^{3}}\frac{g^{4}}{8\pi^{4}}\frac{1}{k^{+}}\int_{0}^{k^{+}}\frac{dq^{+}}{2\pi}\int_{x,y,z,x',w'}e^{-ik\cdot(w'-z)}\notag\\&\quad\times
\frac{e^{-ik\cdot(z-x)\frac{k^{+}-q^{+}}{k^{+}}}}{\big[(k^{+}-q^{+})(y-x)^{2}+q^{+}(y-z)^{2}\big]\,(x'-w')^{2}(x-z)^{2}}\,G_{q}\notag\\&\quad\times\left\{\Big[\frac{k^{+}}{q^{+}}\Big]_{+}-1\right\}\notag\\
&\qquad\times\Big[U^{ab'}(x')-U^{ab'}(w')\Big]\Big[-f^{adc}f^{fbe}U^{db}(x)U^{ce}(z)+N_{c}U^{af}(y)\Big]\rho^{b'}(x')\rho^{f}(y).
\label{eq:finfin}
\end{align}
Here the $q^{+}$ integral has been written as $\int_{0}^{k^{+}}$, using $\int_{k^{+}}^{0}dq^{+}\,T^{im}A^{im}(\cdots)=\int_{0}^{k^{+}}dq^{+}\,\frac{G_{q}}{a^{2}r^{2}}(\cdots)$.

\textbf{Longitudinal logarithms} (these go to the JIMWLK file):
\begin{align}
\frac{d^{3}N_{2\,\log}}{d^{3}k}&=\frac{1}{(2\pi)^{3}}\frac{g^{4}}{16\pi^{5}}\frac{1}{k^{+}}\,{\color{blue}\Big[\log\frac{\vee}{\Lambda}-\log\frac{\vee}{k^{+}}\Big]}\int_{x,y,z,x',w'}\Bigg\{\notag\\&\quad\Bigg[e^{-ik\cdot(w'-z)}\,\frac{(x'-w')\cdot(y-z)}{(x'-w')^{2}(x-z)^{2}(y-z)^{2}}\Big[1+\frac{(y-x)\cdot(x-z)}{2(y-x)^{2}}\Big]\notag\\
&\quad+e^{-ik\cdot(w'-x)}\,\frac{(x'-w')\cdot(y-x)}{(x'-w')^{2}(x-z)^{2}(y-x)^{2}}\Big[1-\frac{(y-z)\cdot(x-z)}{2(y-z)^{2}}\Big]\Bigg]\notag\\&\quad\times\Big[U^{ab'}(x')-U^{ab'}(w')\Big]\Big[-f^{adc}f^{fbe}U^{db}(x)U^{ce}(z)+N_{c}U^{af}(y)\Big]\rho^{b'}(x')\rho^{f}(y)\notag\\
&\quad-2\,e^{-ik\cdot(w'-z)}\,\frac{(x'-w')\cdot(y-z)}{(x'-w')^{2}(y-z)^{2}}\,\frac{(y-x)\cdot(y-z)}{(y-x)^{2}(x-z)^{2}}\notag\\&\qquad\times\Big[U^{ab'}(x')-U^{ab'}(w')\Big]\Big[N_{c}U^{af}(y)-N_{c}U^{af}(z)\Big]\rho^{b'}(x')\rho^{f}(y)\Bigg\}.
\label{eq:finlog}
\end{align}

\section{Checks and remarks}
\begin{enumerate}
\item \textbf{Bookkeeping.} \eqref{eq:finUV}+\eqref{eq:finfin}+\eqref{eq:finlog} reproduce $N_{2a}+N_{2b}$ up to $O(\Lambda/k^{+})$, and nothing is counted twice. The subtraction appears as $+\int S$ in \eqref{eq:finUV} and as $-\int S$ in the last lines of \eqref{eq:finfin} and \eqref{eq:finlog}.
\item \textbf{Only one $\log\mu$.} It is the one in \eqref{eq:finUV}. The finite part and the log part are UV finite in $d=2$:
\begin{itemize}
\item the $\mathcal{C}_{2}\to\mathcal{C}^{(b)}_{z}$ part of their $1/(x-z)^{2}$ singularity is cancelled by the subtraction terms (at every $p^{+}$, and $+1+1-2=0$ for the logs);
\item the $\mathcal{C}^{(a)}$ part vanishes at $x=z$ by colour.
\end{itemize}
\item \textbf{Equivalent compact form.} The blue $2\log\frac{k^{+}}{\Lambda}$ in \eqref{eq:finUV} and the last line of \eqref{eq:finlog} are $\pm2\log\frac{k^{+}}{\Lambda}\int_{x}K_{S}\,\mathcal{C}^{(b)}_{z}$: the same $d$-dimensional integral with opposite signs. They cancel identically. One may therefore drop both, leaving
\begin{itemize}
\item only $-\frac{11}{6}$ in the UV bracket (this is the form used in the Divergent file);
\item only the first two terms in $N_{2\,\log}$.
\end{itemize}
In that form, however, the $\log$ terms must be kept in $d$ dimensions, because separately they diverge at $x\to z$ for the $\mathcal{C}^{(b)}$ colour.
\item \textbf{Where each piece goes.}
\begin{itemize}
\item $-\frac{11}{6}\times$UV goes to the Divergent file, where it combines with Row II's $\frac{11}{3}$.
\item \eqref{eq:finfin} goes to the finite (Other\,terms) file.
\item \eqref{eq:finlog} goes to the JIMWLK file.
\end{itemize}
\item \textbf{Large $x$ (IR, not UV).} At $p^{+}=0$ the first log term falls only like $1/x^{2}$ at large $x$. This large-distance behaviour has to be checked in the sum of all $\log(k^{+}/\Lambda)$ terms of the JIMWLK file.
\end{enumerate}

\end{document}
